\chapter{Interfaccia Utente}

L'intero sistema \`e controllabile tramite una dashboard web realizzata con
\textbf{Streamlit}, avviabile con il comando \texttt{streamlit run app.py}.
L'interfaccia \`e organizzata in una barra laterale di configurazione globale e
due tab principali: \textit{Simulazione} e \textit{Benchmark}.

\section{Barra laterale}

La sidebar consente di selezionare il \textbf{file istanza} (mappa JSON) e il
\textbf{seed} di riproducibilit\`a (valore $-1$ per seed casuale).
Questi parametri sono condivisi da entrambe le modalit\`a.

\section{Tab Simulazione}

Il pannello di simulazione \`e suddiviso in tre colonne:

\begin{itemize}
  \item \textbf{Configurazione} (sinistra) --- permette di impostare il numero di agenti
        (1--10), il limite di tick (100--750) e il delay tra frame.
        Per ogni agente si sceglie la strategia di esplorazione, il raggio di visione
        ($v_r \in \{1,2,3\}$) e il raggio di comunicazione ($c_r \in \{1,2\}$).
        \`E possibile caricare un preset JSON salvato in precedenza o
        personalizzare le icone degli agenti e dei pacchi.
  \item \textbf{Visualizzazione} (centro) --- mostra un rendering in tempo reale della
        griglia, generato frame per frame tramite Pygame in modalit\`a off-screen.
        Ogni cella viene colorata in base al tipo (muro, corridoio, magazzino, ingresso, uscita);
        gli agenti sono cerchi colorati con etichetta ID; i pacchi e i magazzini hanno icone
        dedicate.
        Una barra di avanzamento indica il tick corrente e lo stato della simulazione.
  \item \textbf{Stato} (destra) --- riporta in tempo reale il tick attuale, il conteggio
        degli oggetti consegnati e la batteria residua di ogni agente.
\end{itemize}

Al termine della simulazione vengono mostrati i risultati finali: oggetti consegnati,
tasso di completamento, tick totali, energia media e tempo CPU.
Una tabella riassuntiva per agente elenca strategia, raggio, passi effettuati e
batteria finale. Un grafico mostra la curva cumulativa delle consegne nel tempo.

La simulazione pu\`o essere ripetuta pi\`u volte senza ricaricare la pagina;
uno \textbf{storico} delle run precedenti consente di confrontare le configurazioni
testate nella stessa sessione.
Al termine \`e possibile scaricare il preset corrente in formato JSON per riutilizzarlo.

\section{Tab Benchmark}

La modalit\`a benchmark consente di eseguire automaticamente un numero arbitrario di
simulazioni con configurazioni diverse, per esplorare lo spazio dei parametri.

La configurazione avviene tramite:
\begin{itemize}
  \item \textbf{Strategie}: selezione multipla tra le strategie disponibili oppure
        strategia singola fissa;
  \item \textbf{Raggi}: range di visione e comunicazione, ciascuno in modalit\`a
        casuale (intervallo) o fissa (valore singolo);
  \item \textbf{Numero di preset}: il sistema calcola e mostra lo spazio di ricerca
        totale (combinazioni per agente elevate al numero di agenti) e permette di
        scegliere quanti preset testare.
\end{itemize}

I preset vengono generati senza duplicati (campionamento casuale oppure enumerazione
esaustiva se il numero richiesto copre l'intero spazio).
Durante l'esecuzione una barra di avanzamento mostra il preset corrente e la percentuale
completata.

Al termine vengono prodotti:
\begin{itemize}
  \item Una \textbf{tabella ordinabile} con tutti i risultati (tasso di completamento,
        tick, energia, configurazione, strategia dominante);
  \item Un \textbf{grafico a barre} dei tick di completamento, colorato per strategia
        dominante, con linea della media;
  \item Un \textbf{grafico a barre} degli oggetti consegnati per preset;
  \item Un file \textbf{ZIP esportabile} contenente: \texttt{summary.csv},
        \texttt{results.json} (configurazioni dettagliate),
        \texttt{curves.json} (curve di consegna), \texttt{metadata.json}
        (parametri della sessione) e le immagini dei grafici in formato PNG.
\end{itemize}

Questo workflow ha permesso di generare i dati presentati nel Capitolo~\ref{ch:risultati}.
